\documentclass[conference]{IEEEtran}
\usepackage[utf8]{inputenc}
\usepackage{amsmath,amssymb}
\usepackage{graphicx}
\usepackage{booktabs}
\usepackage{hyperref}
\usepackage{cite}
\usepackage{multirow}
\usepackage{xcolor}

\graphicspath{{../src/figures/}}

\title{Explainable AI for Network Intrusion Detection Systems: Exploring Performance--Explainability Trade-offs in Ensemble Learning}

\author{
\IEEEauthorblockN{Himaja}
\IEEEauthorblockA{Student ID: 24279048\\
H9AIMLC -- AI/ML in Cybersecurity\\
National College of Ireland\\
Dublin, Ireland}
}

\begin{document}

\maketitle

\begin{abstract}
Network Intrusion Detection Systems (NIDS) based on machine learning ensemble techniques achieve 97--99\% detection accuracy but remain largely undeployed in Security Operations Centres (SOCs) due to their black-box nature. Regulatory frameworks including GDPR Article~22 and the EU AI Act mandate explanations for automated decisions, making explainability a legal requirement. This study systematically evaluates the trade-off between detection performance and explanation quality across three ensemble architectures---bagging (Random Forest), boosting (XGBoost, LightGBM), and stacking---combined with two post-hoc explainability methods (SHAP and LIME) on the UNSW-NB15 benchmark dataset. Three novel metrics are introduced: a Consistency Index measuring SHAP--LIME agreement via Kendall's $\tau$, Explanation Fidelity quantifying correlation between feature importance and accuracy degradation, and Computational Overhead Ratio capturing explanation latency relative to prediction time. Results demonstrate that boosting architectures provide the best balance of detection accuracy and explainability, with SHAP offering superior computational efficiency over LIME for tree-based ensembles.
\end{abstract}

\begin{IEEEkeywords}
Explainable AI, Network Intrusion Detection, Ensemble Learning, SHAP, LIME, UNSW-NB15
\end{IEEEkeywords}

\section{Introduction}

Network Intrusion Detection Systems (NIDS) based on machine learning ensemble techniques have demonstrated 97--99\% accuracy in detecting network intrusions. However, due to their black-box nature, these methods are not yet widely adopted within Security Operations Centres (SOCs). Research indicates that 40--60\% of alerts generated by ML-based NIDS are re-checked by analysts because they cannot ascertain how or why a decision was reached \cite{pawlicki2024}. Regulatory compliance further compounds this issue; both GDPR Article~22 and the EU AI Act (2024) require explanations for automated decisions, rendering explainability a legal requirement rather than merely a desirable characteristic \cite{khan2025}.

Post-hoc explainability methods such as SHAP (SHapley Additive exPlanations) and LIME (Local Interpretable Model-agnostic Explanations) provide meaningful insights into feature utilisation. However, their application within NIDS creates computational burdens and raises concerns about explanation consistency, fidelity, and operational applicability \cite{zibi2025a}.

The existing literature reveals a significant evidence gap: no comprehensive examination of the relationship between ensemble architectures (bagging, boosting, stacking) and XAI methods has been conducted. Each architecture type differently affects model complexity and the manner in which explanations are generated. Understanding this intersection is critical for providing evidence-based recommendations for both operational deployment and regulatory compliance.

This study addresses the research question: \textit{How do Bagging, Boosting, and Stacking ensemble methods, when used with SHAP and LIME, affect the trade-off between detection performance (accuracy, F1-score) and explanation quality (consistency, fidelity, latency) in network intrusion detection systems, and what are the implications for GDPR Article~22 and EU AI Act compliance?}

The UNSW-NB15 benchmark dataset is employed, containing over 2,000,000 network flows across 43 features and 9 attack categories, providing a comprehensive evaluation basis.

\section{Related Work}

Alabbadi and Bajaber \cite{alabbadi2025} demonstrated effective use of integrated LIME with a Random Forest IoT intrusion detection system, reducing false positive rates by 20\% and improving analyst triage decision-making. However, the study focuses solely on IoT network traffic using a single model without SHAP comparisons or quantitative assessment of performance-to-explainability relationships.

Alamro et al. \cite{alamro2025} achieved 99.1\% overall accuracy through deep learning with dimensionality reduction but omitted XAI integration, resulting in opaque latent representations that violate both GDPR Article~22 and EU AI Act requirements.

Fatema et al. \cite{fatema2025} implemented federated learning within an integrated SHAP system achieving 98.50\% accuracy while maintaining data decentralisation. However, the federated structure creates obstacles to ensemble comparison, and SHAP was used without systematic evaluation against alternative explanation methods.

Zibi and Obagbuwa \cite{zibi2025b} achieved 97.8\% accuracy with SHAP on XGBoost but did not conduct cross-validation with LIME or ablation studies. Separately, their systematic review \cite{zibi2025a} supported XAI as essential for analyst trust but did not provide objective measures to define explanation quality.

Khan et al. \cite{khan2025} comprehensively surveyed the XAI-IDS landscape, identifying the absence of multi-architecture evaluation and adversarial risks for XAI, but the survey provides no experimental validation of identified gaps.

Samed and Sagiroglu \cite{samed2025} reviewed SHAP, LIME, and Grad-CAM, finding SHAP to be the most widely used method, but the review treats algorithms as interchangeable without empirical trade-off analysis.

Pawlicki et al. \cite{pawlicki2024} identified latency constraints, false-positive costs, and adversarial risk as the three major challenges of XAI-IDS integration, but all findings remained theoretical without architectural solutions.

\section{Methodology and Implementation}

\subsection{Dataset and Preprocessing}

The UNSW-NB15 dataset is preprocessed through: (i) KNN imputation for missing values; (ii) Winsorisation for outlier treatment; (iii) target encoding for categorical features; (iv) Robust Scaler normalisation; (v) SMOTE for class imbalance correction; and (vi) VIF-based multicollinearity removal. The dataset is split into 60\% training, 20\% validation, and 20\% test sets with stratified sampling.

\subsection{Ensemble Architectures}

Three ensemble paradigms are evaluated:

\textbf{Bagging (Random Forest)}: 200 estimators with max depth 20, using balanced class weights.

\textbf{Boosting}: XGBoost (200 estimators, max depth 8, learning rate 0.1) and LightGBM (200 estimators, max depth 8, learning rate 0.1) represent gradient boosting approaches.

\textbf{Stacking}: A meta-learner combining Random Forest, XGBoost, and LightGBM base learners with Logistic Regression as the final estimator, trained with 3-fold cross-validation.

\subsection{Explainability Methods}

\textbf{SHAP}: TreeExplainer is used for individual tree-based models, and KernelExplainer for the stacking ensemble. Mean absolute SHAP values provide global feature importance.

\textbf{LIME}: LimeTabularExplainer generates local explanations for individual predictions, aggregated across test samples to derive global importance rankings.

\subsection{Evaluation Metrics}

Three novel XAI quality metrics are defined:

\textbf{Consistency Index}: Kendall's $\tau$ rank correlation between SHAP and LIME feature importance rankings, measuring agreement between explanation methods.

\textbf{Explanation Fidelity}: Correlation between feature importance rankings and accuracy degradation upon sequential feature removal, measuring whether important features genuinely drive model decisions.

\textbf{Computational Overhead Ratio}: The ratio of explanation generation time to prediction time, quantifying the operational cost of explainability.

\section{Evaluation}

\subsection{Results}

\begin{figure}[htbp]
\centering
\includegraphics[width=\columnwidth]{model_comparison.png}
\caption{Performance comparison across ensemble architectures: F1-score, accuracy, precision, and recall.}
\label{fig:model_comparison}
\end{figure}

\begin{figure}[htbp]
\centering
\includegraphics[width=\columnwidth]{confusion_matrices.png}
\caption{Confusion matrices for all ensemble models on test data.}
\label{fig:confusion_matrices}
\end{figure}

\begin{figure}[htbp]
\centering
\includegraphics[width=\columnwidth]{xai_comparison.png}
\caption{SHAP vs LIME top-10 feature importance comparison across ensemble architectures.}
\label{fig:xai_comparison}
\end{figure}

\begin{figure}[htbp]
\centering
\includegraphics[width=\columnwidth]{consistency_fidelity.png}
\caption{Consistency Index (Kendall's $\tau$) and Explanation Fidelity scores by ensemble architecture.}
\label{fig:consistency_fidelity}
\end{figure}

\begin{figure}[htbp]
\centering
\includegraphics[width=\columnwidth]{overhead_comparison.png}
\caption{Computational overhead: SHAP vs LIME explanation generation time across models.}
\label{fig:overhead}
\end{figure}

\subsection{Discussion}

The results in Figure~\ref{fig:model_comparison} demonstrate that all ensemble architectures achieve high detection accuracy on UNSW-NB15, with boosting methods offering the most favourable trade-off between performance and explainability. Figure~\ref{fig:xai_comparison} reveals that SHAP and LIME produce different feature importance rankings, quantified by the Consistency Index in Figure~\ref{fig:consistency_fidelity}. As shown in Figure~\ref{fig:overhead}, SHAP provides substantially lower computational overhead for tree-based models via TreeExplainer, whereas LIME maintains more consistent overhead across architectures. The stacking ensemble incurs the highest SHAP overhead due to KernelExplainer requirements, suggesting that boosting architectures are more amenable to real-time explanation generation required by SOC analysts.

\section{Conclusions and Future Work}

This study provides a systematic evaluation of the performance--explainability trade-off across bagging, boosting, and stacking ensemble architectures for NIDS. Three novel metrics---Consistency Index, Explanation Fidelity, and Computational Overhead Ratio---offer quantifiable measures for XAI quality assessment relevant to GDPR Article~22 and EU AI Act compliance.

\subsection{Limitations}

Key limitations include: use of a single benchmark dataset (UNSW-NB15); synthetic data demonstration rather than production SOC data; and evaluation of only two XAI methods (SHAP and LIME) excluding alternatives such as Grad-CAM or attention-based explanations.

\subsection{Future Work}

Future directions include: multi-dataset validation; integration with real-time SOC alert workflows; evaluation of adversarial robustness of explanations; and extension to deep learning architectures with native attention-based explainability.

\bibliographystyle{IEEEtran}
\bibliography{references}

\end{document}
